%!TEX encoding = IsoLatin

%
% Chapitre "Introduction"
%

\chapter{Introduction}
\label{s:intro}

Grand Nord est un projet de d�veloppement d'un concept de capteur autonome fixe dans le but de compter ainsi que de documenter la faune arctique. Le Minist�re de la Faune Arctique nous sollicite donc pour construire ce syst�me autonome qui devra �chantillonner les activit�s des lemmings qui se trouvent sur le site arctique en recueillant, entre autres, des images et d'autres donn�es importantes de qualit�s dans le but de produire diverses statistiques. Le syst�me doit �tre robuste et fiable afin d'assurer une interaction externe minimale avec celui-ci advenant le cas d'un probl�me. Il devra aussi �tre en mesure de fonctionner de mani�re autonome pour une dur�e d'une ann�e. Finalement, il est important de prendre en consid�ration les co�ts li�s � la production de ce syst�me autonome et de pr�voir une installation facile ne moyennant pas de comp�tences en ing�nierie.

\bigskip
Ce rapport sera divis� en 6 chapitres afin de bien s�parer chaque concept concernant la mise en oeuvre de ce syst�me.
\begin{itemize}
   \item Description;
   \item Besoins et objectifs;
   \item Cahier des charges;
   \item Conceptualisation et analyse de faisabilit�;
   \item �tude pr�liminaire;
   \item Concept retenu.
\end{itemize}




